\documentclass[conference]{IEEEtran}
\IEEEoverridecommandlockouts
% The preceding line is only needed to identify funding in the first footnote. If that is unneeded, please comment it out.
\usepackage{cite}
\usepackage{amsmath,amssymb,amsfonts}
\usepackage{algorithmic}
\usepackage{graphicx}
\usepackage{textcomp}
\usepackage{xcolor}
\def\BibTeX{{\rm B\kern-.05em{\sc i\kern-.025em b}\kern-.08em
    T\kern-.1667em\lower.7ex\hbox{E}\kern-.125emX}}
\begin{document}

\title{AANEC-Relaxed Trapped Surface Theorem and Stellar Mass-Gap Derivation}\author{Anonymous Research Plan Author}\maketitle

\begin{abstract}
This proposal outlines a formal theoretical framework for classical general relativity, specifically targeting spherically symmetric stellar collapse with isotropic pressure. The planned contribution is the derivation of an AANEC-relaxed trapped-surface theorem. This approach replaces pointwise energy condition assumptions with the Averaged Null Energy Condition (AANEC) combined with an explicit averaged null-convergence/focusing assumption on null congruence generators. The framework aims to establish that reaching a model-dependent compactness threshold forces trapped-surface formation, thereby excluding stable non-black-hole remnants above a specific mass scale within this restricted class.

The expected outcome of this theoretical work is the identification of rigorous boundaries for the proposed theorem. It is anticipated that the analysis will determine whether the added averaged focusing condition is sufficient to yield a trapped-surface bound, or if the theorem must be restricted due to violations of AANEC, incomplete null generators, or the dominance of nonspherical dynamics. This formal result is expected to serve as a consistency constraint for astrophysical mass-gap interpretations, distinguishing geometric exclusion from alternative stellar-physics mechanisms such as explosion energetics or equation-of-state uncertainties.
\end{abstract}
\begin{IEEEkeywords}
AANEC, trapped surfaces, stellar collapse, mass gap, compactness threshold, null energy condition, regular matter fields
\end{IEEEkeywords}\section{Introduction}
Recent gravitational-wave observations have inferred population properties of compact binary mergers, categorizing them into binary black hole, binary neutron star, and neutron star-black hole classes. These studies provide a foundation for examining the mass distributions of stellar remnants.~\cite{cite_fd0f77b2a173c0ab}

While models suggest that the location of the lower edge of the black hole mass gap is dictated by near-final core mass, alternative stellar-physics mechanisms such as explosion energetics, fallback, and equation-of-state uncertainties also influence remnant formation. This creates a need for a formal geometric framework that can be clearly separated from these astrophysical confounders.~\cite{cite_7a407c865faafd11}

The proposed research design assumes spherically symmetric collapse of massive stars with isotropic pressure, global hyperbolicity, and regular matter fields. It further assumes that matter fields satisfy the Averaged Null Energy Condition (AANEC) along complete null generators, and that an explicit averaged null-convergence/focusing condition holds.

This proposal plans to develop an AANEC-relaxed trapped surface theorem and derive a compactness bound for non-black-hole stellar remnants. The contribution will focus on replacing pointwise energy condition assumptions with AANEC plus an explicit averaged focusing assumption, aiming to exclude stable non-black-hole remnants above a model-dependent compactness threshold within this restricted class.

\section{Background, Survey, and Research Gap}
Classical singularity theorems rely on pointwise energy conditions to guarantee the focusing of null geodesics and the formation of trapped surfaces. However, quantum field theory indicates that pointwise energy conditions are systematically violated by realistic matter fields, necessitating the use of weaker averaged statements such as the Achronal Averaged Null Energy Condition (AANEC). While AANEC offers a robust pathway to preserve singularity theorems in the face of quantum violations, the specific assumptions and validity domains under which Penrose-type theorems support black hole existence remain under-specified in the current evidence plan.~\cite{cite_0ed84243cf171bc4}

Astrophysical observations of stellar-mass black holes reveal a distinct paucity of compact objects in the 2-5 solar mass range, known as the lower mass gap. This gap is widely interpreted as a discriminator for supernova engine physics, favoring rapid explosion dynamics over delayed fallback scenarios. The location of the lower edge of this gap is dictated by the near-final core mass, which sets the resulting black hole mass, and appears robust against variations in metallicity and stellar wind mass loss.~\cite{cite_7a407c865faafd11}

Despite the observational consensus regarding the mass gap, the identification of a precise boundary variable that separates black hole and non-black-hole outcomes remains unresolved in the relevant evidence cluster for theoretical stellar evolution. Current models struggle to definitively isolate a single variable that accounts for the transition, particularly when accounting for uncertainties in explosion and unbinding energies across different progenitor masses.

\section{Research Questions and Planned Contributions}
This section outlines the central research questions and planned theoretical contributions for the proposed mechanism replacement. The primary objective is to formulate a qualified trapped-surface bound for non-black-hole stellar remnants under relaxed energy conditions, while explicitly separating formal geometric results from astrophysical interpretations of the stellar mass gap.

\begin{itemize}
\item The following questions guide the formal development and verification plan:
\item What is the precise mathematical form of the explicit averaged null-convergence condition (A5) that distinguishes it from standard AANEC (A2)?
\item How is the Misner-Sharp compactness parameter C formally defined for the specific metric class (D4) under consideration?
\item How is the stability of non-black-hole remnants (A6) formally defined to distinguish them from regular black hole models like Hayward or Bardeen?
\item Do candidate witnesses CE1 and CE2 genuinely fail condition A5, or is A5 currently ill-defined in this context?
\item Should numerical verification of the Raychaudhuri equation and trapped surface formation be conducted for a discrete subset of test metrics?
\end{itemize}

\begin{itemize}
\item The planned contributions include:
\item Proposing a formal theorem for classical GR spherical collapse that replaces pointwise energy conditions with AANEC plus an explicit averaged null-convergence/focusing assumption.
\item Deriving a trapped-surface threshold via Raychaudhuri-type focusing and Misner-Sharp compactness within a restricted model class.
\item Bounding stable non-black-hole remnants through TOV-like stellar-structure assumptions.
\item Treating the observed 2-5 solar mass gap as an astrophysical consistency constraint rather than a direct theorem-derived prediction, thereby acknowledging alternative stellar-physics mechanisms.
\end{itemize}

The design assumes spherically symmetric collapse with isotropic pressure, global hyperbolicity, regular matter fields, and complete null generators. It further assumes that matter fields satisfy AANEC without uncontrolled quantum ANEC violations. These assumptions define the scope of the formal claims and are treated as bounded design constraints pending human review of specific mathematical definitions.

\section{Idea Source Checkpoints and Direction Selection Audit}
The selected direction, mechanism\_replacement, is documented across three available audit snapshots. These checkpoints originate from idea\_candidate.json, idea\_portfolio.json, and idea\_result.json. The temporal order of these snapshots is unknown, as no iteration IDs, timestamps, or auditable diffs are provided to establish chronology. Consequently, these records are treated as neutral audit snapshots rather than evidence of iterative refinement or temporal evolution.

The proposed research design assumes a formal theorem/conjecture for classical general relativity, specifically focusing on spherical nonrotating collapse with isotropic pressure. The scope includes regular matter fields, global hyperbolicity, complete null generators, AANEC, and an explicit averaged null-convergence/focusing condition. This design explicitly excludes quantitative derivations of the observed 2-5 solar mass gap and does not claim a general proof that all massive stars necessarily form black holes, treating the mass gap instead as an astrophysical consistency constraint.

\section{Problem Definition, Assumptions, and Hypotheses}
\paragraph{Definition.} The formal framework defines the mass of the collapsing matter or remnant object (M), the areal radius or radial coordinate (r), the expansion scalar of a null geodesic congruence (theta), the stress-energy tensor (T), the spacetime metric (g), and the tangent vector to null geodesic generators (k). The Misner-Sharp compactness parameter (C) is proposed as a central variable but requires human input to establish its specific mathematical formulation and effective mass definition within the chosen metric.

\begin{itemize}
\item The theoretical model relies on the following declared assumptions: (1) Spacetime and matter fields satisfy the conditions of classical General Relativity with spherically symmetric collapse of massive stars and isotropic pressure; (2) Matter fields satisfy the Averaged Null Energy Condition (AANEC) along complete null generators; (3) Spacetime is globally hyperbolic and sufficiently regular for trapped-surface arguments; (4) Null generators in the relevant region are complete; (5) An explicit averaged null-convergence/focusing condition holds along the relevant null generators; (6) Non-black-hole remnants are modeled as stable equilibrium objects obeying a bounded compactness or TOV-like limit.
\end{itemize}

\paragraph{Proposition (proposed).} Two candidate formal propositions are proposed. Proposition P1 posits that in spherically symmetric stellar collapse satisfying AANEC, global hyperbolicity, regular matter fields, complete null generators, and an explicit averaged null-convergence/focusing condition, reaching a model-dependent compactness threshold forces trapped-surface formation. Proposition P2 posits that stable non-black-hole remnants above a compactness-dependent mass scale are excluded within the restricted model class defined by AANEC and averaged focusing.

The research hypotheses map these propositions to testable formal structures. Hypothesis H1 links the focusing and compactness threshold to trapped surface formation, with a decision rule that it is considered proved only if proof obligations PO1 and PO2 are closed for all admissible metrics in the restricted model class. Hypothesis H2 links the exclusion of stable non-black-hole remnants above a maximum compactness to the derivation steps S3-S4, with a decision rule that it is considered proved only if those steps follow from the stable equilibrium assumption and the compactness definition; otherwise, it is rejected if a witness satisfying all assumptions is found.

The study design assumes a mathematical/theoretical approach where the experimental unit is a formal proof and counterexample search. It is assumed that the observed 2-5 solar mass gap is not a direct consequence of the theorem but is treated as an astrophysical consistency constraint subject to alternative stellar-physics explanations. The design further assumes that the added averaged null-convergence/focusing condition is an independent geometric hypothesis, not a consequence of AANEC alone.

\section{Study Design and Methods}
The study employs a mathematical and theoretical design focused on formal proof and counterexample search. The primary method involves replacing pointwise energy condition assumptions with the Averaged Null Energy Condition (AANEC) combined with an explicit averaged null-convergence/focusing condition. The analysis will proceed by deriving trapped-surface thresholds via Raychaudhuri-type focusing equations and Misner-Sharp compactness parameters, followed by bounding stable non-black-hole remnants through TOV-like stellar-structure assumptions. This design treats the observed stellar mass gap as an astrophysical consistency constraint rather than a direct theorem prediction.

\begin{itemize}
\item The theoretical framework operates under the following bounded design assumptions: (1) spherical symmetry and isotropic pressure for the collapse model; (2) global hyperbolicity and sufficient regularity for trapped-surface arguments; (3) complete null generators; (4) matter fields satisfying AANEC without uncontrolled quantum ANEC violations; and (5) non-black-hole remnants modeled as stable equilibrium objects obeying bounded compactness limits.
\end{itemize}

\paragraph{Human-review checklist.} Several methodological components require human input to resolve undefined operational parameters and governance protocols. Specifically, the mathematical form of the explicit averaged null-convergence condition and the definition of the Misner-Sharp compactness parameter must be finalized. Additionally, the measurement plan, quality control procedures, and data governance strategies remain undefined and require specification before the formal proof or simulation phases can proceed.

\section{Expected Outcome Branches and Conditional Conclusions}
\paragraph{Conditional outcome branch.} The expected outcomes are conditional and unobserved. If the prespecified analysis is consistent with the declared relation while planned controls do not favor a stated alternative explanation, the result would support, but not prove, the declared relation within the design boundary. If the prespecified analysis indicates variation across declared conditions, units, or measurement contexts, the relation may be conditional or heterogeneous; no universal conclusion is warranted. If the prespecified comparison does not support the declared relation or instead favors a declared alternative explanation, the proposed relation is not supported in this design boundary; absence of support is not proof of absence generally. If prespecified quality-control, missingness, protocol-deviation, or validity criteria prevent interpretation, no scientific conclusion is warranted because the planned design did not yield interpretable evidence.

Conditional conclusions are restricted to the declared boundary conditions. The theorem fails or must be restricted if AANEC is violated, the added averaged null-convergence/focusing condition fails, null generators are incomplete, quantum stress-energy violates ANEC, nonspherical/rotating/magnetic dynamics dominate, or stable exotic compact objects evade the averaged-focusing assumptions. The mass-gap inference fails or becomes non-discriminating if explosion energetics, fallback, equation-of-state physics, pair-instability processes, or binary evolution dominate the observed remnant distribution.

\section{Risks, Limitations, and Human Review Requirements}
The proposed formal reasoning plan requires qualified human review before proceeding, as the current status of the formal derivation and counterexample analysis is unresolved. Specifically, the mathematical definitions for the Misner-Sharp compactness parameter and the explicit averaged null-convergence/focusing condition are missing, preventing precise verification of the theorem's assumptions and boundaries. Additionally, the distinction between formal counterexamples and empirical alternatives in the high-compactness regime requires expert evaluation to ensure the validity of the proposed witnesses.

\section{Definitions, Propositions, and Proof Obligations}
\paragraph{Definition.} D1: Mass of the collapsing matter or remnant object.

\paragraph{Definition.} D2: Areal radius or radial coordinate.

\paragraph{Definition.} D3: Expansion scalar of a null geodesic congruence.

\paragraph{Definition.} D4: Misner-Sharp compactness parameter. The specific formal definition, including the formula for effective mass within radius r in the chosen metric, requires human input.

\paragraph{Definition.} D5: Stress-energy tensor.

\paragraph{Definition.} D6: Spacetime metric.

\paragraph{Definition.} D7: Tangent vector to null geodesic generators.

\paragraph{Proposition (proposed).} P1: In spherically symmetric stellar collapse satisfying AANEC, global hyperbolicity, regular matter fields, complete null generators, and an explicit averaged null-convergence/focusing condition, reaching a model-dependent compactness threshold forces trapped-surface formation.

\paragraph{Proposition (proposed).} P2: Stable non-black-hole remnants above a compactness-dependent mass scale are excluded within the restricted model class defined by AANEC and averaged focusing.

PO1: The explicit mathematical form of the averaged null-convergence condition (A5) distinguishing it from A2 is not provided, preventing verification of whether candidate witnesses satisfy the condition.

PO2: This proof obligation remains unresolved.

PO3: This proof obligation remains unresolved.

\section{Forward Derivation and Counterexample Search Plan}
The proposed forward derivation outlines a path from initial assumptions regarding spherically symmetric collapse and energy conditions to the exclusion of stable non-black-hole remnants above a compactness threshold. This plan relies on a sequence of proposed steps: first, applying a Raychaudhuri-type focusing equation to demonstrate that the expansion scalar becomes negative along null generators under the explicit averaged focusing condition (S1). Second, identifying that if the compactness parameter exceeds a model-dependent threshold, both outgoing and ingoing null expansions become negative, indicating trapped surface formation (S2). Third, establishing that trapped surface formation is incompatible with the definition of a stable non-black-hole equilibrium object below the maximum compactness bound (S3). Finally, deriving the exclusion bound for stable remnants (S4). It is emphasized that these steps are currently unverified formal proposals. The derivation depends on the resolution of missing formal definitions, specifically the Misner-Sharp compactness parameter (D4) and the mathematical form of the averaged null-convergence condition (A5/PO1), which remain subjects for human input.

The counterexample search plan targets the validity of the proposed theorem by attempting to construct witnesses that satisfy all stated assumptions while negating the conclusion. Two candidate classes are considered: hypothetical static, spherically symmetric, regular metrics with high compactness that evade trapped surfaces, and regular black hole models (e.g., Hayward or Bardeen) that might allow for remnant interpretations. The plan involves checking each candidate against assumptions A1 through A6. Preliminary analysis suggests that candidates may fail to satisfy the explicit averaged focusing condition (A5) or the stability bound (A6), or may be classified as black holes rather than non-black-hole remnants due to the presence of horizons. The validity of these counterexamples remains unverified and is contingent on the precise formal definitions of A5 and D4. The search method is currently limited to LLM-proposed candidates, requiring rigorous mathematical verification to determine if they constitute valid counterexamples or merely boundary cases.

Several limitations constrain the current formal reasoning plan. The absence of a formal mathematical definition for the Misner-Sharp compactness parameter (D4) prevents precise calculation of thresholds and rigorous verification of the derivation steps. Similarly, the lack of a specific mathematical form for the averaged null-convergence condition (A5) makes it impossible to verify whether candidate counterexamples actually satisfy this assumption. The distinction between formal counterexamples and empirical alternatives requires human review of the specific mathematical structures proposed. Furthermore, the current status of the proof obligations (PO1, PO2, PO3) is unresolved, indicating that the logical chain from assumptions to conclusion has not been established. These gaps necessitate human input to define the missing formal elements before the theorem can be considered verified or falsified.

\appendices
\section{Idea Source Checkpoints and Direction Selection Audit}
This section presents the available audit snapshots for the selected research direction, 'mechanism\_replacement'. These checkpoints are derived from the idea generation and selection pipeline. They are not presented as a temporal evolution or iteration history, as the chronological order of these snapshots is unknown.

\begin{itemize}
\item The following checkpoints are available:
\item Checkpoint 1: Sourced from 'idea\_candidate.json' (root). Contains the central hypothesis regarding AANEC-relaxed trapped surface theorem and stellar mass-gap derivation.
\item Checkpoint 2: Sourced from 'idea\_portfolio.json' (/selected\_primary\_idea). Contains a qualified version of the AANEC trapped-surface bound for non-black-hole stellar remnants.
\item Checkpoint 3: Sourced from 'idea\_result.json' (/directions/0). Contains the detailed hypothesis mapping, including assumptions, boundary conditions, and evidence requirements for the mechanism replacement direction.
\end{itemize}

\section{Variables, Symbols, and Operational Definitions}
\paragraph{Definition.} The formal reasoning framework introduces the following candidate definitions: D1 (mass of the collapsing matter or remnant object, symbol M), D2 (areal radius or radial coordinate, symbol r), D3 (expansion scalar of a null geodesic congruence, symbol theta), D5 (stress-energy tensor, symbol T), D6 (spacetime metric, symbol g), and D7 (tangent vector to null geodesic generators, symbol k). The Misner-Sharp compactness parameter (D4, symbol C) requires further human input to establish its rigorous mathematical formulation.

The operationalization of specific physical variables is currently incomplete. V6 requires the formal definition of the Misner-Sharp compactness parameter. V7 requires a specified measurement method for stellar remnant mass. V8 requires a defined metric for explosion energy. V10 requires the specification of equation of state models or nuclear physics parameters. V11 requires the measurement or initial condition specification for rotation. V12 requires the measurement or initial condition specification for magnetic fields.

\section{Evidence Coverage, Unknown Items, and Review Checklist}
The evidence base for this proposal is limited to registered sources and evidence cards. The primary support for relaxing pointwise energy conditions to averaged forms comes from the finding that quantum fields systematically violate pointwise conditions, necessitating weaker averaged statements for broader validity.~\cite{cite_0ed84243cf171bc4}

\begin{itemize}
\item Key limitations identified in the current design include: (1) The mathematical form of the explicit averaged null-convergence/focusing condition (A5) is missing, preventing precise verification of whether proposed counterexamples satisfy this premise. (2) The definition of Misner-Sharp compactness (D4) is absent, hindering the calculation of specific thresholds for stable remnant exclusion. These gaps require formal specification before the theorem can be rigorously tested.
\end{itemize}

\paragraph{Human-review checklist.} Before proceeding with the formal derivation and counterexample search, the following review steps are required: (1) Human review must define the mathematical form of assumption A5 (explicit averaged null-convergence/focusing). (2) Human review must supply the formal definition of D4 (Misner-Sharp compactness parameter). (3) The counterexample analysis status must be upgraded from 'design\_assumption' to 'verified' or 'rejected' only after these definitions are resolved and the witnesses are checked against them.\section*{References}
\nocite{cite_07d6df5d6978a1fc,cite_0ed84243cf171bc4,cite_1e54e3c2b35110ad,cite_2576833bf7ef143a,cite_283d00782f1f9ecb,cite_357dd780a9e3a378,cite_364386a3569b6e66,cite_37ecbded9d629353,cite_42e341ff6471ce89,cite_572fec0a1e5d74d9,cite_5ea4dc0e1b288757,cite_63476e0b002da386,cite_72d7ff391b29c495,cite_789b47fb0241d9ae,cite_7a407c865faafd11,cite_7a44e6730c44cb39,cite_7ae92f06684c60ef,cite_977568cf7b8a55bf,cite_98c5e596c0c34b68,cite_bfae3a92efd1d48f,cite_c369df652181e8f7,cite_cbbced417410bc6c,cite_cc5bc07a219047b4,cite_d3016ce6d1faa8ec,cite_e228b5fd31e37646,cite_e7b42f305bc08498,cite_ecea60a2300e4892,cite_f634820e22d7a238,cite_fd0f77b2a173c0ab}
\bibliographystyle{IEEEtran}
\bibliography{references}
\end{document}
